\begin{textsample}{POS Dim 4 – 2017\_c6 – Score 11.00 – 2017\_c6/enade\_2017\_eng\_prod\_cf6e5e5935.txt}  \label{ex:f4_pos_033}
CAPÍTULO 6 ANÁLISE TÉCNICA DA PROVA Este capítulo tem por objetivo apresentar o desempenho dos estudantes concluintes de Engenharia de Produção no Enade/2017.
Para isso, foram calculadas as estatísticas básicas da prova em seu todo ( seção 6.1.1 ), bem como as estatísticas relacionadas aos Componentes de Formação Geral ( seção 6.1.2 ) e de Conhecimento Específico da Área ( seção 6.1.3 ).
Dadas as suas características, foram analisadas, em separado, as questões objetivas ( seção 6.2 ) e as questões discursivas ( seção 6.3 ).
Para as questões objetivas, foram disponibilizados os índices de facilidade e de discriminação ponto-bisserial, também em separado, para os Componentes de Formação Geral ( seção 6.2.1 ) e de Conhecimento Específico ( seção 6.2.2 ).
De cada componente, uma das questões foi escolhida para exemplificar a análise gráfica, relacionando as alternativas escolhidas pelos estudantes ( inclusive o gabarito ) com o número de acertos no componente.
O Anexo I apresenta a íntegra da análise gráfica para todas as questões objetivas.
Para cada uma das questões discursivas, os conteúdos dos tipos mais comuns de respostas dos estudantes são apresentados e comparados com o padrão de respostas esperado ( ver Anexo VIII com o padrão de respostas ).
Tomando como base as duas questões discursivas do Componente de Formação Geral, a seção 6.3.1.6 apresenta comentários sobre a correção das respostas com respeito à Língua Portuguesa.
As tabelas apresentam as seguintes estatísticas das notas19 : média do desempenho na prova, erro padrão da média, desvio padrão, nota mínima, mediana e nota máxima para cada um de seus componentes.
Tais estatísticas contemplam o total de estudantes concluintes da Área de Engenharia de Produção inscritos e presentes à prova do Enade/2017, tendo em vista agregações, ou por Grandes Regiões e o país como um todo, ou por Categoria Administrativa, Organização Acadêmica da IES e Modalidade de Ensino.
Em relação aos gráficos de distribuição de notas, o intervalo considerado foi de 10 unidades, aberto à esquerda e fechado à direita, com exceção do primeiro intervalo, [ 0 ; 10 ], fechado em ambos os extremos.
Para os gráficos de distribuição das notas das questões discursivas, foram consideradas mais duas categorias : questão em branco20 e nota zero. 19 Uma definição dessas estatísticas pode ser encontrada no Glossário. 20 Nesse grupo estão incluídas também as respostas classificadas como nulas ou desconsideradas. 6.1 ESTATÍSTICAS BÁSICAS DA PROVA Esta seção apresenta estatísticas selecionadas e histogramas da nota geral ( 6.1.1 ) e de cada componente : Formação Geral ( 6.1.2 ) e Conhecimento Específico ( 6.1.3 ).
São também apresentadas estatísticas selecionadas de subpopulações caracterizadas por Grande Região, Categoria Administrativa, Organização Acadêmica e Modalidade de Ensino. 6.1.1 Estatísticas Básicas Gerais A Tabela 6.1 apresenta as Estatísticas Básicas da prova, por Grande Região, dos estudantes concluintes de Engenharia de Produção.
A Média das notas da prova como um todo ( nas seções seguintes serão analisados os Componentes de Formação Geral e de Conhecimento Específico ) foi 41,7, sendo que os alunos da região Norte obtiveram a Média mais baixa ( 37,8 ), e os da região Sul obtiveram a Média mais alta ( 44,1 ).
As demais médias foram : 40,8 na região Nordeste, 41,6 na região Sudeste e 41,1 na região Centro-Oeste.
O Desvio padrão para o Brasil como um todo foi 12,7, sendo o maior Desvio padrão encontrado na região Sudeste ( 12,9 ) e o menor, na região Norte ( 11,0 ), indicando uma dispersão um pouco menor das notas desta última região.
A região que obteve a maior nota Máxima foi a Sul ( 88,5 ), ao passo que a região que atingiu a menor nota Máxima foi a Norte ( 74,9 ).
A Mediana do Brasil como um todo foi 40,9, sendo a maior Mediana obtida na região Sul ( 43,5 ), e a menor obtida na Norte ( 37,3 ).
A nota Mínima foi zero na região Sudeste, 4,0 na região Nordeste, 6,4 na região Sul, 8,4 na região Centro-Oeste e 10,0 na região Norte.
Considerando -se as notas segundo Grande Região, observa -se que existe diferença estatisticamente significativa ao nível de 95 \% entre a menor Média, obtida na região Norte ( 37,8 ), e as médias de todas as demais regiões²¹. ²¹ Todas as comparações deste capítulo utilizam os intervalos de 95 \%.
Os erros-padrão da média que possibilitam os testes estão disponíveis nas tabelas.
O comportamento das notas dos estudantes de todo o Brasil pode ser observado no Gráfico 6.1 que apresenta um histograma com a distribuição das mesmas.
Essa é uma distribuição unimodal com a moda no intervalo ( 30 ; 40 ], seguido pelo intervalo ( 40 ; 50 ].
A Tabela 6.2 apresenta informações referentes à Média da nota final desagregadas por Categoria Administrativa, por Organização Acadêmica e por Modalidade de Ensino.
A Média das notas da prova como um todo foi 41,7.
Em relação à Categoria Administrativa, os estudantes das IES Públicas obtiveram Média mais alta ( 49,0 ), e os das IES Privadas obtiveram Média mais baixa ( 39,4 ), que a Média nacional.
Observa -se que existe diferença estatisticamente significativa entre as médias das notas das IES Públicas e Privadas.
A diferença entre as médias das regiões Norte e Sul ( 6,3 ), a maior e a menor Média, é inferior à diferença entre IES Públicas e Privadas ( 9,6 ), caracterizando uma maior diversidade administrativa do que regional.
O Desvio padrão para as IES Públicas ( 13,2 ) foi superior ao do Brasil como um todo ( 12,7 ), indicando uma dispersão um pouco maior das notas nesta Categoria Administrativa.
No tocante à Organização Acadêmica, as Universidades e os CEFET/IFET obtiveram Média mais alta que a nacional ( 44,0 e 47,6 respectivamente ).
As Médias das Faculdades e a dos Centros Universitários foram menores que a nacional, 38,1 e 39,8, respectivamente.
Constata -se que existe diferença estatisticamente significativa ao nível de 95 \% entre a Média mais baixa, 38,1 das Faculdades, e a dos demais tipos de Organização Acadêmica.
A Média da Modalidade Educação a Distância ( 37,1 ) foi inferior à Média da Educação Presencial ( 41,9 ), e a diferença entre elas é estatisticamente significativa ao nível de 95 \%. 6.1.2 Estatísticas Básicas no Componente de Formação Geral A Tabela 6.3 apresenta as Estatísticas Básicas em relação ao componente da prova que avalia a Formação Geral dos estudantes concluintes.
Os alunos de todo o Brasil obtiveram desempenho médio de 54,6.
Quanto à variabilidade, o Desvio padrão das notas dos estudantes do Brasil, como um todo, foi 16,9.
A maior Média foi obtida na região Sul ( 56,3 ), e a menor, na região Norte ( 52,1 ).
As demais médias foram : 53,8 na região Nordeste, 54,6 na região Sudeste e 54,5 na região Centro-Oeste.
Já o maior Desvio padrão foi obtido nas regiões Nordeste e Centro-Oeste ( 17,1 ), e o menor, na região Norte ( 15,9 ).
Os demais desvios padrões foram : 17,0 na região Sudeste e 16,4 na região Sul.
A maior nota no Componente de Formação Geral da prova do Enade/2017 foi 98,8, obtida por, pelo menos, um aluno na região Sul.
A menor nota Máxima foi obtida na região Centro-Oeste ( 95,0 ).
A Mediana do Brasil, como um todo, foi 55,5, sendo a menor Mediana encontrada na região Norte ( 52,5 ), e a maior, encontrada na região Sul ( 57,4 ).
A nota Mínima nesta parte foi zero em todas as regiões, com exceção da região Norte, onde a nota Mínima foi 7,5.
Considerando -se as notas segundo Grande Região, observa -se que existe diferença estatisticamente significativa entre a maior Média das notas do Componente de Formação Geral, obtida na região Sul ( 56,3 ), e as médias obtidas nas demais regiões.
O Gráfico 6.2 propicia a avaliação do desempenho dos estudantes do Componente de Formação Geral, a partir do histograma da distribuição das notas correspondentes.
A distribuição é unimodal, com moda em ( 50 ; 60 ], intervalo modal superior ao da distribuição de notas da prova como um todo, ( 30 ; 40 ] ( Gráfico 6.1 ).
Nota -se, ainda que, no Gráfico 6.2, as notas apresentam dispersão maior do que a do Gráfico 6.1 ( distribuição das notas da prova ), confirmado pela comparação dos desvios padrões : 12,7 para a nota da prova como um todo e 16,9 para o Componente de Formação Geral.
Na Tabela 6.4 são apresentadas as informações referentes ao desempenho dos concluintes do Componente de Formação Geral, em diferentes agregações : Categoria Administrativa, Organização Acadêmica e Modalidade de Ensino.
Observa -se que existe diferença estatisticamente significativa entre as médias dos tipos de Categoria Administrativa : a maior Média obtida por estudantes de IES Públicas ( 62,0 ) e a menor pelos de IES Privadas ( 52,3 ).
Nota -se que também há diferença estatisticamente significativa entre todas as médias dos quatro tipos de Organização Acadêmica, valor maior para CEFET/IFET ( 61,9 ) e menor para Faculdades ( 51,1 ).
No que se refere à Modalidade de Ensino, a Educação Presencial apresentou maior Média do que a Educação a Distância ( médias 54,8 e 50,9, respectivamente ).
Tais médias apresentam diferenças estatisticamente significativas. 6.1.3 Estatísticas Básicas do Componente de Conhecimento Específico A Tabela 6.5 apresenta as Estatísticas Básicas referentes ao Componente de Conhecimento Específico da área de Engenharia de Produção.
A Média do desempenho dos alunos do Brasil, como um todo, foi 37,4.
A maior Média foi obtida na região Sul ( 39,9 ), e a menor, na região Norte ( 33,0 ).
As demais médias foram : 36,5 na região Nordeste, 37,2 na região Sudeste e 36,7 na região Centro-Oeste.
Quanto à variabilidade das notas, o Desvio padrão do Brasil, como um todo, foi 13,8, sendo o maior Desvio padrão observado na região Sudeste ( 13,9 ), e o menor, na região Norte ( 12,2 ).
Os demais desvios foram : 13,4 na região Nordeste, 13,3 na região Sul e 13,7 na região Centro-Oeste.
A Mediana das notas dos estudantes de todo o Brasil foi 36,4.
A maior Mediana ocorreu na região Sul ( 39,8 ), e a menor, na região Norte ( 32,4 ).
As demais medianas foram : 35,3 na região Nordeste e 36,3 nas regiões Sudeste e Centro-Oeste.
A nota Máxima do Brasil, como um todo, foi 91,3, sendo obtida por, pelo menos, um aluno na região Sudeste.
As demais notas máximas foram : 72,0 na região Norte, 84,3 na região Nordeste, 91,2 na região Sul, e 85,0 na região Centro-Oeste.
A nota Mínima foi zero nas regiões Sudeste e Sul.
Já nas demais regiões a nota Mínima foi maior do que zero : 4,1 na região Norte, 5,3 na região Nordeste e 2,0 na região Centro-Oeste.
Observa -se que existe diferença estatisticamente significativa da menor Média das notas do Componente de Conhecimento Específico, obtida na região Norte ( 33,0 ), em relação à Média das demais regiões.
Assim como os Gráficos 6.1 e 6.2, o Gráfico 6.3, apresentado a seguir, permite uma avaliação do desempenho de concluintes de Engenharia de Produção em relação ao Componente de Conhecimento Específico com um histograma da distribuição das notas correspondentes.
Esta também é uma distribuição unimodal, e o grupo modal é o ( 30 ; 40 ], mesmo grupo modal da prova como um todo e abaixo do grupo modal para a Formação Geral.
A Tabela 6.6 apresenta uma comparação dos resultados em relação à Categoria Administrativa, à Organização Acadêmica e a Modalidade de Ensino, agora levando em conta o desempenho de alunos do Componente de Conhecimento Específico da prova.
No que se refere à Organização Acadêmica, a maior Média foi dos CEFET/IFET ( 42,8 ), vindo a seguir, em ordem decrescente, a das Universidades ( 39,7 ), a dos Centros Universitários ( 35,5 ) e das Faculdades ( 33,8 ).
O maior Desvio padrão foi o das Universidades ( 14,4 ), superior ao do valor para o Brasil ( 13,8 ), como um todo.
As Universidades também obtiveram a maior nota Máxima ( 91,3 ).
Os Centros Universitários obtiveram nota Máxima igual a 83,5, as Faculdades obtiveram nota Máxima igual a 85,8 e os CEFET/IFET igual a 83,2.
As medianas foram : 38,8 nas Universidades, 34,8 nos Centros Universitários, 32,7 nas Faculdades e 42,5 nos CEFET/IFET, a maior.
A nota Mínima foi diferente de zero apenas no grupo de estudantes dos CEFET/IFET.
Observa -se que existem diferenças estatisticamente significativas ao nível de 95 \% no Componente de Conhecimento Específico entre as notas das Organizações Acadêmicas.
Quanto à Categoria Administrativa, observa -se um comportamento semelhante àquele da parte de Formação Geral e à prova como um todo, ou seja, existe diferença estatisticamente significativa entre as médias das IES Públicas ( 44,6 ), a maior, e das IES Privadas ( 35,1 ).
Também semelhante é o comportamento das médias no que toca à Modalidade de Ensino, há diferença estatisticamente significativa entre a Média da Educação Presencial ( 37,6 ) e a da Educação a Distância ( 32,5 ). 6.2 ANÁLISE DAS QUESTÕES OBJETIVAS Esta seção apresenta estatísticas selecionadas e histogramas das Questões Objetivas de Formação Geral ( 6.2.1 ) e de Conhecimento Específico ( 6.2.2 ).
São também apresentadas e comparadas as médias das subpopulações caracterizadas por Grande Região. 6.2.1 Componente de Formação Geral A Tabela 6.7 apresenta as Estatísticas Básicas relativas às oito questões objetivas do componente da prova que abrange a Formação Geral dos estudantes.
A Média do Brasil foi 53,4.
A menor Média foi encontrada na região Norte ( 49,5 ), e a maior, na região Sul ( 55,9 ).
As demais médias foram : 52,4 na região Nordeste e 53,2 na região Sudeste e 53,1 na região Centro-Oeste.
O Desvio padrão do Brasil foi 21,2, sendo o maior Desvio padrão encontrado na região Centro-Oeste ( 21,5 ), e o menor, na região Norte ( 20,1 ).
Os demais desvios foram : 21,3 nas regiões Nordeste e Sudeste, e 20,6 na região Sul.
As medianas do Brasil como um todo e de quatro das cinco regiões foi 50,0 – a única região com valor diferente de Mediana foi a Sul, com 62,5.
A nota Máxima 100,0, assim como a nota Mínima ( 0,0 ) foram iguais para todas as regiões.
A Tabela 6.8 apresenta o Índice de Facilidade e o Índice de Discriminação ( ponto-bisserial ) para cada uma das questões objetivas do Componente de Formação Geral.
Quanto ao Índice de Facilidade, foram usadas as seguintes cores para diferenciar o nível de dificuldade da questão :  Azul para as questões classificadas com índice Muito fácil ( \&gt ; =0,86 ), verde para as questões classificadas com índice Fácil ( 0,61 a 0,85 ), amarelo para as questões classificadas com Médio ( 0,41 a 0,60 ), vermelho para as questões classificadas com Difícil ( 0,16 a 0,40 ) e roxo para as questões classificadas com Muito difícil ( \&lt ; =0,15 ).
Já quanto ao Índice de Discriminação, foram usadas as seguintes cores para qualificar a questão : • As questões classificadas com índice Fraco receberam a cor vermelho ( \&lt ; =0,19 ), as classificadas com Médio receberam a cor amarelo ( 0,20 a 0,29 ), as classificadas com Bom receberam a cor verde ( 0,30 a 0,39 ) e as classificadas com Muito bom ( \&gt ; =0,40 ) receberam a cor azul.
As questões objetivas do Componente de Formação Geral, segundo o Índice de Facilidade, foram assim avaliadas : das oito questões, uma teve o Índice de Facilidade classificado como Muito fácil, por ter índice de acertos 0,87, e duas questões foram tidas como Fácil.
Três questões foram consideradas com índice de dificuldade Médio, situando -se no intervalo entre 0,41 e 0,60 do Índice de Facilidade, ou seja, houve entre 41,0 \% e 60,0 \% de acertos, enquanto duas questões foram classificadas na categoria Difícil, situando -se no intervalo entre 0,16 e 0,40.
Por fim, nenhuma questão apresentou menos de 15 \% de acertos, razão pela qual não houve questão classificada como Muito difícil.
O Índice de Facilidade variou de 0,28 a 0,87, e o de Discriminação, de 0,40 a 0,51.
A Tabela 6.9 apresenta a distribuição das questões considerando -se simultaneamente a classificação dos dois índices.
Todas as questões objetivas de Formação Geral foram classificadas com Índice de Discriminação Muito bom e figuraram entre os níveis de dificuldade de Muito fácil a Difícil.
A questão 6, classificada na categoria Muito fácil do Índice de Facilidade, teve Índice de Discriminação 0,43.
Na categoria Fácil, as questões 4 e 5, tiveram Índice de Discriminação 0,51 e 0,48, respectivamente.
As três da categoria Médio, do índice de Facilidade, ( questões 2, 3, e 8 ) também tiveram Índice de Discriminação variando de 0,40 a 0,46.
Por fim, duas questões consideradas na categoria Difícil, questões 1 e 7, tiveram poder de discriminação 0,47 e 0,45, respectivamente.
Em particular, a questão 4 foi a que apresentou o maior poder discriminatório, com índice 0,51, e foi considerada Fácil em termos de facilidade, com uma proporção de 0,62 acertos.
O máximo de acertos foi alcançado pela questão 6 com um Índice de Facilidade de 0,87.
A questão 1 foi considerada a mais difícil, com índice de facilidade 0,28.
O Gráfico 6.4, para exemplificar, analisa o comportamento da questão de número 4 de Formação Geral.
Trata -se de uma questão considerada Fácil, com relação à facilidade e a que obteve o maior Índice de Discriminação dessa parte da prova.
Neste gráfico, cada uma das cinco curvas representa o percentual de respostas em determinada alternativa da questão, em função do número de acertos dos estudantes nessa parte da prova ( Formação Geral/Múltipla Escolha ), antes de possíveis eliminações pelo critério do ponto-bisserial.
A curva em verde corresponde à alternativa B, a alternativa correta para esta questão.
Observa -se que entre os estudantes com menor número de acertos, nessa parte do exame, a situação mais frequente foi a escolha da alternativa E ( em vermelho ), incorreta.
Por exemplo, entre os estudantes que acertaram três questões, 26,2 \% escolheram a alternativa E, 42,9 \% escolheram a alternativa B ( em verde ), a alternativa correta, 18,6 \% escolheram a alternativa D ( em roxo ), 7,3 \% escolheram a alternativa A ( azul ) e 4,8 \% a C ( preto ).
Entre os que acertaram três respostas entre as questões de múltipla escolha de Formação Geral, 0,1 \% deixou a questão em branco e 0,1 \% marcou mais de uma alternativa, invalidando a questão.
À medida em que o número de acertos aumenta, indicando desempenho melhor nessa parte da prova, aumenta concomitantemente a proporção de estudantes que selecionaram a alternativa correta B, atingindo 100 \% para os estudantes com 8 acertos.
Essa análise permite verificar como a questão discriminou os grupos de desempenho, justificando o alto índice obtido na questão.
Cumpre notar que não é possível inferir deste gráfico nem o índice de facilidade ( que seria uma média da proporção ponderada pela quantidade de alunos com cada uma das notas ), nem o índice de discriminação ponto-bisserial, por razão equivalente.
No caso extremo no qual a grande concentração dos acertos dos alunos fosse abaixo de 4, o índice de facilidade seria obrigatoriamente abaixo de 40 \% ( neste exemplo ).
Caso a concentração fosse em 6 acertos ou mais, o índice seria obrigatoriamente acima de 60 \%.
Os gráficos relativos às demais questões de Formação Geral constam do Anexo I. 6.2.2 Componente de Conhecimento Específico A Tabela 6.10 apresenta as Estatísticas Básicas em relação às questões objetivas do Componente de Conhecimento Específico da prova de Engenharia de Produção por Grande Região.
Nessa parte da prova uma questão foi anulada pela Comissão Assessora da Área, questão 35.
Além disso, como será discutido ainda nessa seção, quatro questões objetivas do Componente de Conhecimentos Específicos não foram usadas no cômputo das notas por terem sido descartadas pelo critério do ponto-bisserial.
Assim, as notas foram calculadas com base em 22 das 27 questões objetivas de Conhecimentos Específicos.
A Média do Brasil deste componente foi de 41,7.
A menor Média foi observada na região Norte ( 36,8 ), e a maior, na região Sul ( 44,5 ).
O Desvio padrão de todo o Brasil foi 15,6, sendo o menor Desvio padrão encontrado na região Norte ( 14,1 ), e o maior, na região Sudeste ( 15,8 ).
A Tabela 6.11 apresenta os Índices de Facilidade e Discriminação ( ponto-bisserial ) das questões objetivas do Componente de Conhecimento Específico para os estudantes de Engenharia de Produção.
Para facilitar a diferenciação das questões, usaremos as mesmas cores da Tabela 6.8 para as diferentes classificações dos Índices de Facilidade e de Discriminação.
Observa -se que uma questão foi anulada pela CAA e as análises dizem respeito às 26 questões válidas.
A partir do Índice de Facilidade obtido, pode -se concluir que mais de metade das questões objetivas da prova foi considerada, pelo menos, Difícil : das 26 questões válidas, 16 foram classificadas como Difícil ou como Muito difícil.
Uma questão foi classificada como Fácil, e outras nove consideradas como Médio.
Já quanto aos índices de discriminação das questões objetivas do Componente de Conhecimento Específico da prova, tem -se como resultado a seguinte classificação : 12 das 26 questões foram consideradas como boas, enquanto apenas duas tiveram Índice de Discriminação Muito bom.
Assim, para um pouco mais de metade das questões válidas – 14 em 26 – os Índices de Discriminação foram Bom ou Muito bom.
Dentre as demais, oito delas foram classificadas como Médio e as outras quatro como Fraco, sendo 12, por conseguinte, a quantidade de questões nos dois patamares mais baixos de discriminação.
Constata -se, assim, que a prova – no que se refere ao Componente de Conhecimento Específico – possuía capacidade média de discriminar entre aqueles que dominavam ou não o conteúdo.
A Mediana de todo o Brasil foi 40,9, a mesma das regiões Nordeste, Sudeste e Centro-Oeste.
A Mediana foi maior que a nacional na região Sul ( 45,5 ) e menor na região Norte ( 36,4 ).
A nota Máxima da prova foi 100,0, obtida nas questões objetivas do Componente de Conhecimento Específico, por, pelo menos, um aluno da região Sudeste.
A nota Mínima foi zero nas regiões Sudeste, Sul e Centro-Oeste e igual a 4,5 nas regiões Norte e Nordeste.
A Tabela 6.11 apresenta os Índices de Facilidade e Discriminação ( ponto-bisserial ) das questões objetivas do Componente de Conhecimento Específico para os estudantes de Engenharia de Produção.
Para facilitar a diferenciação das questões, usaremos as mesmas cores da Tabela 6.8 para as diferentes classificações dos Índices de Facilidade e de Discriminação.
Observa -se que uma questão foi anulada pela CAA e as análises dizem respeito às 26 questões válidas.
A partir do Índice de Facilidade obtido, pode -se concluir que mais de metade das questões objetivas da prova foi considerada, pelo menos, Difícil : das 26 questões válidas, 16 foram classificadas como Difícil ou como Muito difícil.
Uma questão foi classificada como Fácil, e outras nove consideradas como Médio.
Já quanto aos índices de discriminação das questões objetivas do Componente de Conhecimento Específico da prova, tem -se como resultado a seguinte classificação : 12 das 26 questões foram consideradas como boas, enquanto apenas duas tiveram Índice de Discriminação Muito bom.
Assim, para um pouco mais de metade das questões válidas – 14 em 26 – os Índices de Discriminação foram Bom ou Muito bom.
Dentre as demais, oito delas foram classificadas como Médio e as outras quatro como Fraco, sendo 12, por conseguinte, a quantidade de questões nos dois patamares mais baixos de discriminação.
Constata -se, assim, que a prova – no que se refere ao Componente de Conhecimento Específico – possuía capacidade média de discriminar entre aqueles que dominavam ou não o conteúdo.
A Tabela 6.12 apresenta a distribuição das questões considerando -se simultaneamente a classificação dos dois índices.
Dentre as questões que alcançaram os maiores índices de discriminação, classificadas com Índice de Discriminação Muito bom, uma delas, a de número 22, foi a índice mais elevado ( 0,42 ), foi classificada como Difícil quanto ao Índice de Facilidade, com 23 \% dos estudantes marcando a opção correta.
A questão de número 12 foi a mais difícil dentre as 26 questões específicas válidas, com baixo Índice de Facilidade, apenas 14 \% de acertos.
Essa questão apresentou poder discriminatório muito baixo, 0,10, o que comprova ter sido a mais difícil para os estudantes.
Destaca -se, também, a questão 11, com Índice de Facilidade 0,33, o que, em termos percentuais, corresponde a 33 \% de estudantes que responderam acertadamente.
Já 0,07 foi o seu de Índice de Discriminação.
Tais questões foram, portanto, pelo critério ponto-bisserial, consideradas inadequadas.
Por isso, as questões 11 e 12, além das questões 17 e 24 foram eliminadas do cômputo da nota final.
A título de exemplo das análises do comportamento das questões objetivas, o Gráfico 6.5 analisa a questão 22 do Componente de Conhecimento Específico.
Esta questão foi considerada pelas respostas dos estudantes avaliados na prova como uma questão com nível de facilidade Difícil, 0,23, ou seja, 23 \% dos estudantes assinalaram acertadamente a opção D, correspondente ao gabarito.
Como já comentado, seu Índice de Discriminação foi igual a 0,42, classificado como Muito bom, o maior valor de discriminação.
Neste gráfico, cada uma das cinco curvas representa o percentual de respostas em determinada alternativa da questão 22, em função do número de acertos dos estudantes nessa parte da prova, antes de possíveis eliminações de questões pelo critério do ponto-bisserial.
A alternativa correta D, representada no gráfico pela curva em roxo, foi escolhida em maiores proporções pelos alunos com desempenho melhor nessa parte da prova.
Já as alternativas incorretas, também denominadas distratores, foram selecionadas, principalmente, por aqueles com notas mais baixas.
Observa -se que a soma não é 100 \% por causa das questões não respondidas ou com mais de uma opção marcada.
Aqueles com nota zero, na sua quase totalidade deixaram esta questão em branco ou marcaram mais de uma alternativa, comportamento considerado inválido.
A proporção de alunos que selecionou a resposta correta D aumenta gradativamente, chegando a atingir 100 \% apenas para os que obtiveram 22 acertos.
A proporção dos que escolheram alternativa incorreta E ( em vermelho ) só é superada pela proporção alternativa correta a partir de 13 acertos, como função do número de acertos nessa parte da prova.
Os gráficos relativos às demais questões do Conhecimento Específico constam do Anexo I. 6.3 ANÁLISE DAS QUESTÕES DISCURSIVAS Esta seção apresenta estatísticas selecionadas e histogramas das Questões Discursivas de Formação Geral ( 6.3.1 ) e Conhecimento Específico ( 6.3.2 ).
São também apresentadas e comparadas as médias de subpopulações caracterizadas por Grande Região. 6.3.1 Componente de Formação Geral As análises dos resultados de desempenho dos estudantes de Engenharia de Produção nas duas questões discursivas relativas à Formação Geral encontram -se na Tabela 6.13 e no Gráfico 6.6.
Na Tabela 6.13, observa -se que a nota Média nesse conjunto de questões foi acima da obtida nas objetivas.
Os estudantes, de todo o Brasil, obtiveram, em Formação Geral, Média 53,4 nas questões objetivas e 56,5 nas questões discursivas.
Pode -se notar também que o Desvio padrão nesse conjunto de questões foi um pouco maior do que o obtido nas objetivas, 21,2 nas questões objetivas e 23,0 nas questões discursivas.
A maior Média foi obtida na região Sul ( 56,9 ), e a menor, na região Nordeste ( 55,8 ).
A Mediana de todo o Brasil, neste componente, foi 62,0, a mesma sendo obtida nas regiões Nordeste, Sudeste e Sul.
Na região Centro-Oeste a Mediana foi maior ( 62,5 ) e na região Norte foi menor ( 61,5 ).
A nota Máxima ( 98,0 ) foi a mesma nas regiões Nordeste e Sudeste, sendo 96,0 na região Norte, 97,5 na região Sul e 97,0 na região Centro-Oeste.
A nota Mínima ( 0,0 ) foi a mesma em todas as regiões do Brasil.
O Gráfico 6.6 representa a distribuição das notas nas questões discursivas do Componente de Formação Geral.
A moda desta distribuição ocorre no intervalo ( 60 ; 70 ].
Destaca -se também o intervalo [ 0 ; 10 ] com distribuição próxima a 8 \% do total de notas, sendo que no intervalo [ 0 ; 10 ] incluem -se além da nota zero, a frequência de alunos que deixaram as respostas a este tipo de questão em branco.
Na sequência, os resultados verificados para cada uma das questões discursivas de Formação Geral serão apresentados, estabelecendo -se relações com os conteúdos abordados em cada uma delas.
Os comentários da Banca de docentes corretores, a respeito do observado na correção das respostas dos estudantes, suas impressões e conclusões serão apresentados junto à análise de cada questão.
Cumpre esclarecer que, tendo em vista que as questões discursivas de Formação Geral são padronizadas, ou seja, constam de todas as provas, os comentários da Banca são os mesmos para todas as carreiras acadêmicas, sendo direcionados a todos os estudantes que participaram do Enade/2017.
A seguir, serão analisados os desempenhos da Área de Engenharia de Produção nas duas questões discursivas de Formação Geral do Enade/2017, comparando -se os resultados obtidos com comentários para cada questão. 6.3.1.1 Análise de Conteúdo da Questão Discursiva 1 do Componente de Formação Geral Os dados de Engenharia de Produção, obtidos a partir das respostas à questão 1, encontram -se na Tabela 6.14 e no Gráfico 6.7.
Nessa questão – de melhor desempenho dentre as duas de Formação Geral – os alunos, de todo o Brasil, tiveram Média 57,7.
A maior Média para a questão 1 foi obtida na região Norte ( 59,2 ), e a menor, na região Sudeste ( 57,6 ).
Quanto à variabilidade das notas, o Desvio padrão de todo o Brasil foi 26,8.
O menor Desvio padrão foi obtido na região Norte ( 26,2 ), e o maior Desvio padrão foi obtido na região Nordeste ( 27,4 ).
As medianas das regiões Norte, Nordeste e Sudeste, bem como do Brasil como um todo, foram iguais ( 65,0 ).
Nas regiões Sul e Centro-Oeste a Mediana foi 60,0.
As notas máximas e mínimas da questão discursiva 1 foram, sem exceção, as mesmas para todas as regiões do Brasil, respectivamente, 100,0 e 0,0.
O Gráfico 6.7 mostra a distribuição das notas na questão discursiva 1 do Componente de Formação Geral.
Observa -se que a moda corresponde aos alunos que obtiveram notas no intervalo ( 70 ; 80 ], com aproximadamente 24 \% do total.
Em torno de 7 \% dos participantes deixaram a resposta a essa questão em branco. 6.3.1.2 Comentários sobre a correção de Conteúdo das respostas à Questão Discursiva 1 O enunciado era claro e indicava um tema de extrema relevância, qual seja, a incidência de doença grave integrante do rol das DST ( Doenças Sexualmente Transmissíveis ), que encaminha para a análise do direito fundamental à saúde, que tem \textbf{assento} constitucional ( CRFB, art. 6º. caput : São direitos sociais a educação, a saúde, a alimentação, o trabalho, a moradia, o transporte, o lazer, a segurança, a previdência social, a proteção à maternidade e à infância, a assistência aos desamparados, na forma desta Constituição ).
A questão focalizava, por meio de três textos, um tema social extremamente relevante, qual seja, a relação entre homem e mulher nas suas interações íntimas, no caso representadas pela vida sexual, e o reflexo das relações de poder veiculadas pela sociedade, que apontam, no mais das vezes, o predomínio do homem.
Tal posição de superioridade é justificada pelas convenções arraigadas que acarretam a educação dos homens e mulheres por padrões patriarcais, que redundam em posições machistas.
Um dos textos demandava uma reflexão sobre a realidade dos \textbf{casais} que convivem em relações estáveis ou fugazes e não têm o controle dos efeitos das relações sexuais, que abrangem a gravidez não planejada e, em decorrência do não cuidado, a inserção de doenças transmitidas sexualmente.
Esse panorama é refletido na pouca influência das mulheres sobre os homens quanto à proteção, normalmente representada pelo uso de preservativos nas relações intimas, e após a constatação da doença, no encaminhamento das pessoas para tratamento.
Como efeito colateral, está a transmissão de doenças para os fetos.
Outro texto trazia à tona o maior cuidado que as mulheres têm, em geral, com sua saúde, notadamente no momento em que se descobre grávida e que redunda na realização de exames obrigatórios de pré-natal, que podem facilmente detectar doenças e, caso encontradas, ter o seu tratamento encaminhado.
A par disso, destaca -se que a crise de recursos financeiros agravou a situação, que já recebe influxos sociais negativos, como a dificuldade do homem em utilizar preservativos, realizar exames e procurar tratamentos quando adoece.
Como muito divulgado pelas mídias, a crise financeira do país repercutiu diretamente na prestação de serviços de saúde em todas as esferas : distrital, federal, estadual e municipal.
Assim, agravou -se o problema da falta de medicamentos essenciais para o tratamento da sífilis e, certamente, de outras doenças.
Como se verifica, a temática indicada no enunciado, representada pelos textos, pertence à realidade conhecida amplamente pelos estudantes, pois também é divulgada por todos os meios de comunicação.
Trata -se de tema geral que, no entanto, normalmente, não se inclui nas grades curriculares de cursos superiores que não sejam da área de Saúde.
Apesar disso, a temática é de conhecimento geral e relacionada aos direitos humanos, esses sim, ponto de integração com as diretrizes curriculares.
A linguagem utilizada era clara, indicando a existência dos fatos vinculados a questões relevantes e à saúde pública.
No entanto, como os textos trazem uma ampla gama de possibilidades de interpretação, ocorreram algumas compreensões diversas das esperadas.
Além disso, ressalta -se que a amplitude dos textos disponibilizados no enunciado fornece elementos que poderiam ser simplesmente reproduzidos pelos concluintes nas suas respostas.
O enunciado da questão encaminha a resposta para dois aspectos.
O estudante deveria abordar conteúdo pertinente à vulnerabilidade das mulheres às DST e o papel social do homem em relação à prevenção dessas doenças.
Para responder a esses aspectos, um número expressivo de concluintes utilizou somente os elementos dos textos motivadores, sem nada acrescentar.
A resposta deveria conter também uma segunda parte, a indicação de duas ações voltadas para o público masculino, a serem adotadas no âmbito das políticas públicas de saúde ou de educação, para a redução do problema.
Verificou -se que o desempenho dos estudantes nessa questão foi de médio para bom.
Ao se analisarem as respostas, constata -se que muitos foram os que apresentaram os tópicos presentes no padrão de resposta, amplo o suficiente para albergar as respostas apresentadas pelos concluintes.
Tal resultado pode ser justificado, tendo em vista tratar -se de tema bastante difundido nas diversas mídias, por ter sido apresentado de forma ampla pelos textos motivadores presentes no enunciado.
Por outro lado, o debate no ambiente universitário sobre as relações de poder entre homens e mulheres tende a ser marcante pelo movimento de empoderamento feminino, liderado por organizações sociais com ampla atuação nos estabelecimentos de nível superior.
De forma recorrente, as relações patriarcais entre homens e mulheres foram vinculadas ao modelo machista vigente.
Essas condições, segundo os concluintes, explicam a incidência da doença e a situação de vulnerabilidade das mulheres, bem como a função social dos homens.
Outras respostas apontaram o poder público, nas áreas de saúde e/ou educação, como órgãos responsáveis pela solução do problema.
Para os que se fixaram na área da saúde, foram frequentes a indicação da necessidade de mais medicamentos, vacinas, educação sexual, médicos e postos de atendimentos.
Alguns, no entanto, apenas indicaram ações individuais para o homem, e outros, ações individuais para as mulheres.
Também figuraram respostas que apresentavam ações da igreja como solução, além de outras que faziam referências a campanhas no seio das famílias. 6.3.1.3 Análise de Conteúdo da Questão Discursiva 2 do Componente de Formação Geral A Tabela 6.15 mostra que o desempenho médio dos estudantes na questão discursiva 2 ( Média 55,0 ) foi inferior ao obtido na questão discursiva 1 ( Média 57,7 ).
A região Sul foi aquela cuja Média foi maior ( 55,6 ), e a de menor Média foi a região Norte ( 53,2 ).
Quanto à variabilidade das notas, o Desvio padrão de todo o Brasil foi 31,7, superior ao obtido na questão discursiva 1 ( 26,8 ).
O maior desvio nessa questão foi obtido na região Norte ( 32,5 ), enquanto o menor foi obtido na região Sul ( 30,5 ).
A Mediana de todo o Brasil foi 70,0, a mesma de todas as regiões.
As notas máximas ( 100,0 ) e as notas mínimas ( 0,0 ) também foram as mesmas em todas as regiões, sem exceção.
O Gráfico 6.8 mostra a distribuição das notas na questão discursiva 2 do Componente de Formação Geral.
Observa -se que a maior frequência corresponde aos alunos que obtiveram nota dentro do intervalo ( 70 ; 80 ].
Destacam -se, também, os alunos que deixaram a resposta à questão em branco.
Nota -se, ainda, que as notas ficaram menos concentradas em comparação à questão discursiva de número 1.
Isto pode ser constatado pela comparação do desvio padrão das notas da questão discursiva 2 ( 31,7 ) e da questão discursiva 1 ( 26,8 ). 6.3.1.4 Comentários sobre a correção de Conteúdo das respostas à Questão Discursiva 2 O enunciado era claro e abordava um tema de extrema relevância : a questão de gênero.
Trata -se de uma proposta polêmica, pois abrange a análise da inserção social de minoria vinculada à questão de gênero.
O enunciado solicitava que o concluinte discorresse sobre a importância do nome para as pessoas transgêneras e propusesse uma medida, no âmbito das políticas públicas, para facilitar o acesso dessas pessoas à cidadania.
A respeito da primeira parte, a importância do nome, esperava -se que o estudante abordasse as dificuldades que o uso do nome civil ( nome de registro, nome de batismo, nome da certidão de nascimento, por exemplo ) pode acarretar para as pessoas transgêneras, além de demonstrar conhecimento sobre a importância do nome social, direito já reconhecido através, principalmente, de decisões judiciais e atos administrativos, como decretos e portarias que dispõem sobre o uso do nome social.
Em atendimento à segunda parte do comando, o participante deveria propor uma política pública associada ao direito à cidadania das pessoas transgêneras.
Sobre esse aspecto, esperava -se a compreensão do significado de política pública, ou seja, propostas de âmbito governamental, e não, iniciativas de Organizações Sociais, igrejas ou grupos sociais de qualquer tipo.
O nascimento das pessoas, no Brasil, é objeto de inscrição em \textbf{cartórios} de Registro Civil das Pessoas Naturais, órgão responsável por estabelecer a certificação pública dos dados de todos os indivíduos que não somente nascem, mas casam, eventualmente se divorciam e vêm a óbito.
Daí o nome civil.
Para flexibilizar a correção, outras nomenclaturas foram consideradas, a saber : identificação de nascimento, nome de batismo, nome de registro, nome próprio, nome verdadeiro, que diferenciam o nome civil daquele que o indivíduo deseja utilizar para traduzir a sua opção pessoal de gênero, que se convencionou denominar nome social.
Quanto ao entendimento do termo transgênero, deve -se destacar que o movimento de mudança do gênero é contemporâneo e não está necessariamente vinculado à mudança de sexo.
Existem várias hipóteses em jogo.
Uma é a que envolve cirurgias de mudança de sexo ( a menina extirpa os seus seios e inclui órgão genital por intervenção cirúrgica ; o menino inclui seios e extirpa o seu órgão genital e inclui o órgão feminino também por intervenção cirúrgica ).
Há outro movimento que envolve apenas as vestimentas e o linguajar corporal, além dos relacionamentos ( as meninas assumem o trajar, o andar, os gostos dos meninos e mantêm relacionamentos afetivos com meninas ; os meninos assumem o trajar, o andar, os gostos das meninas e mantêm relacionamentos afetivos com meninos ).
Em relação à discussão da primeira solicitação do comando da questão, destaca -se que todos têm nome civil, que foi escolhido pelos seus pais ou responsáveis e, na falta deles, por autorização do Juiz de Direito, pautado no sexo detectado no nascimento.
Já o nome social é aquele que indica a preferência do indivíduo na sua identificação nas relações sociais e que destoa do seu gênero de registro civil.
Ressalte -se que o Brasil não possui ainda legislação específica sobre pessoas transgêneras.
Observa -se um relativo progresso no reconhecimento de direitos, através, principalmente, de decisões judiciais e atos administrativos, como decretos e portarias que dispõem sobre o uso do nome social.
Sem uma lei que defina os procedimentos da alteração dos documentos para pessoas transexuais, essa parcela da população LGBT ( lésbicas, gays, bissexuais, transexuais e transgêneros ) é obrigada a procurar na Justiça o reconhecimento de sua identidade, em processos que podem ser longos e que dependem do olhar de determinados, profissionais, como médicos, psicólogos, assistentes sociais, bem como do entendimento dos \textbf{advogados}, defensores públicos, juízes e do Judiciário, como um todo.
Atualmente, tramita na Comissão de Direitos Humanos da Câmara dos Deputados o Projeto de Lei João Nery ( 5002/2013 ), dos deputados Jean Willys ( PSOL-RJ ) e Erika Kokay ( PT-DF ), que determina que o reconhecimento da identidade de gênero seja um direito do cidadão.
O projeto recebeu o nome do primeiro trans-homem operado no Brasil.
Grande parte dos respondentes considerou a incongruência entre a imagem corporal, representada pelo gênero assumido pelos trans, e o nome que apresenta nos seus documentos como causa de sérios constrangimentos e embaraços e as dificuldades para que esses sujeitos acessem direitos básicos, como saúde, educação e empregabilidade em uma sociedade heteronormativa e preconceituosa.
Muitos demonstraram reconhecer que a identidade e expressão de gênero, através do nome social, pressupõe que sejam respeitadas a dignidade da pessoa humana, a intimidade, a liberdade e o direito à cidadania plena.
No entanto, também foram encontradas, em quantidades significativas, respostas que demonstravam a existência de uma interpretação social conservadora que não admite a mudança de sexo nem do nome social.
Existe um acentuado preconceito contra todos os movimentos que buscam identificação social diversa da civil.
Isso inclui os variados grupos : gays, travestis, transgêneros, dentre outros.
Já em relação ao segundo tópico a ser abordado pelo concluinte, é notório que as pessoas trans encontram -se destituídas de direitos no campo das políticas públicas, como saúde, habitação, emprego, educação, justiça, dentre outros.
Nesse sentido, uma grande diversidade de medidas, no âmbito das políticas públicas, poderia ser proposta.
Ações voltadas para a capacitação de profissionais da educação, saúde, segurança pública, dentre outras áreas, para melhor atender às demandas e especificidades desse público alvo seriam necessárias.
Ainda em relação à educação, seriam também necessárias ações que contribuíssem para diminuir a evasão, pois o ambiente escolar é difícil para uma pessoa transgênera, por conta das situações discriminatórias e preconceituosas vindas dos colegas de classe, e mesmo das/dos professores e/ou de outros profissionais da área, não podem ser mais um \textbf{complicador} nesse processo.
Ainda no que tange ao assunto, percebe -se que, ao serem excluídas dos bancos escolares, as oportunidades no mercado de trabalho formal praticamente inexistem.
Um outro aspecto que se evidencia no cotidiano desse segmento da sociedade diz respeito às situações de violência sofridas pelas pessoas transgêneras.
Nesse sentido, algumas propostas foram elencadas para minimizar o problema e combater a violência.
Políticas envolvendo melhor formação de funcionários das polícias civil e militar, por exemplo, somariam esforços na qualificação das abordagens policiais e dos atendimentos prestados pelas forças de segurança, pois o respeito à orientação sexual e à identidade e expressão de gênero é, antes de tudo, o respeito à dignidade, à cidadania e à própria democracia.
Nesse sentido, revela -se essencial que, ao sofrerem agressões, sejam elas físicas, psicológicas ou morais no espaço público ou domiciliar, as pessoas trans possam ser atendidas, com respeito inclusive ao nome social, nas delegacias comuns e, inclusive, nas Delegacias Especiais de Atendimento à Mulher ( DEAM ), pois esse é o órgão institucional que possivelmente detém as qualificações necessárias para o atendimento ao gênero feminino ao qual as travestis e mulheres trans efetivamente pertencem.
Também ocorre uma latente inacessibilidade das pessoas trans à assistência jurídica.
Nesse sentido, é necessário implementar projetos de capacitação para profissionais do poder judiciário, como um todo, de modo a prestarem a assistência jurídica gratuita na defesa intransigente dos direitos sociais, civis, trabalhistas e políticos dessa população.
Dessa maneira, a Justiça agiria de forma mais acessível e sensível para o segmento, que terá suas especificidades observadas a começar pelo devido respeito do direito ao uso do nome social.
Nesse âmbito, o mais comum foi propor a adoção de medidas punitivas para quem viola o direito à autodeterminação de gênero e a elaboração de leis que garantam a mudança do nome e assegurem outros direitos para as pessoas trans.
Propostas envolvendo a formação de profissionais, como as mencionadas aqui, e de conscientização da população foram as mais frequentes.
Muitas outras enfatizavam o direito ao uso respeitado do nome social e a garantia de alteração do registro civil, independentemente da cirurgia de redesignação sexual ou transgenitalização.
Muitos questionaram o fato de não bastar a autodeclaração no processo de alteração de nome, enfatizando que o nome social é uma questão de princípio de dignidade e cidadania.
Apesar de pouco mencionada nas respostas corrigidas, políticas envolvendo a promoção da saúde são muito necessárias.
Cabe aos gestores de saúde ações direcionadas, não somente à prevenção e tratamento do HIV/Aids, mas também políticas de saúde que contemplem as especificidades dessa população, tais como o Ministério da Saúde, especialmente o SUS ( Sistema Único de Saúde ), incluir, de fato, o atendimento para pessoas trans para acompanhamento da terapia hormonal, que muitas demandam, respeitando -se o direito de utilização do nome social.
De modo geral, são muitas as políticas públicas a que precisam ser implantadas.
Urge a necessidade de uma agenda mais compromissada por parte dos gestores do nosso país com essa comunidade, que ainda é muito discriminada.
Logo, tais políticas devem contemplar suas necessidades básicas como o direito de acesso aos estudos, à capacitação para o mercado de trabalho, à profissionalização, ao acesso a bens e serviços de qualidade em saúde, habitação, segurança, cultura e assistência social e, certamente, o respeito ao uso do nome social, o que também contribui para essa inclusão.
Assim, as opções apresentadas no padrão de resposta funcionaram como exemplos de políticas públicas possíveis, já que seria difícil apresentar uma lista suficientemente ampla para abranger a grande variedade de respostas encontradas.
O desempenho dos estudantes nessa questão foi médio.
As respostas, de modo geral, indicaram a necessidade de uma atualização por parte dos concluintes quanto a temas essenciais e, no caso da presente questão, a análise sobre aspectos fundamentais da questão de gênero.
Muitos, possivelmente motivados pelos textos apresentados no enunciado, restringiram -se a discorrer sobre preconceito de gênero ou mesmo sobre o preconceito de modo mais geral, sem responder ao que foi solicitado no comando da questão.
Outros apresentaram discursos religiosos ou ideológicos, também sem tratarem do que foi solicitado. 6.3.1.5 Análise de Língua Portuguesa das Questões Discursivas do Componente de Formação Geral Os dados de Engenharia de Produção, obtidos a partir das respostas às questões discursivas do Componente de Formação Geral, no que tange à Língua Portuguesa, encontram -se na Tabela 6.16 e no Gráfico 6.9.
Nesse aspecto, os alunos, de todo o Brasil, obtiveram Média 57,5.
A maior Média com respeito à Língua Portuguesa foi obtida na região Sul ( 58,3 ), e a menor, na região Norte ( 56,1 ).
Quanto à variabilidade das notas, o Desvio padrão de todo o Brasil foi 20,6.
O menor Desvio padrão foi obtido na região Norte ( 19,5 ) e o maior Desvio padrão foi obtido na região Nordeste ( 21,3 ).
A Mediana das notas de Língua Portuguesa foi 62,5 para o Brasil como um todo e para todas as cinco Grandes Regiões.
A nota Máxima foi 95,0, obtida por, pelo menos, um aluno da região Sudeste.
Nas regiões Norte e Centro-Oeste a nota Máxima foi 90,0 e nas regiões Nordeste e Sul foi 92,5.
Já a nota Mínima foi zero em todas as regiões do país.
O Gráfico 6.9 mostra a distribuição das notas de Língua Portuguesa do Componente de Formação Geral.
Observa -se que a maior frequência ( pouco acima de 30 \% ) corresponde aos alunos que obtiveram nota no intervalo ( 60 ; 70 ].
Destaca -se, também, os alunos que deixaram em branco ambas as respostas às questões de Formação Geral, representando em torno de 7 \% do total. 6.3.1.6 Comentários sobre a correção das respostas de Formação Geral com respeito à Língua Portuguesa Os enunciados apresentados em 2017 não explicitaram a exigência de elaboração de um texto “ dissertativo ”, como nos anos anteriores.
O participante deveria fazer uma exposição de seus conhecimentos sobre cada um dos assuntos e estruturar seus textos de acordo com as características do registro formal adequado à situação comunicativa – avaliação de conhecimentos.
Essa configuração determina exigências quanto aos seguintes aspectos : adequação da seleção vocabular, desenvolvimento do conteúdo, organização lógica das ideias, estruturação sintática dos períodos, utilização de procedimentos de encadeamento textual e de referenciação, obediência às exigências morfossintáticas próprias da modalidade escrita da norma-padrão, respeito às regras ortográficas e às regras de acentuação gráfica.
O padrão de resposta utilizado na avaliação das questões 1 e 2 considerou os aspectos relevantes ao bom desempenho linguístico como competências distintas, de modo a permitir um mapeamento detalhado do domínio dos recursos disponíveis na Língua Portuguesa para a comunicação escrita formal.
Com base nesse objetivo, foram avaliados os seguintes aspectos : ( a ) estruturação textual condizente com o gênero solicitado e o modo de organização textual expositivo adequado ao gênero – essa competência envolve : • estruturação sintática condizente com o padrão da modalidade escrita formal da língua portuguesa de modo a garantir a clareza necessária ; • distribuição do conteúdo do texto em parágrafos, de modo a garantir a sua organização temática ; • utilização de operadores discursivos que contribuam para a progressão temática do texto, estabelecendo relações lógicas entre as ideias apresentadas, tanto do ponto de vista intrafrasal, como do interfrasal ; • utilização de procedimentos de referenciação lexical e pronominal que permitam a retomada de referentes textuais ; • utilização de sinais de pontuação que contribuam para a organização lógica da frase e do texto ; • inteligibilidade relacionada ao atendimento das exigências de estruturação textual.
Espera -se, portanto, que o participante recorra a procedimentos linguístico-discursivos para organizar seu texto, permitindo o encadeamento lógico entre suas partes de forma a garantir a progressão e a coerência textuais.
Isso significa que os seguintes procedimentos foram considerados inadequados, de acordo com o padrão de resposta proposto :  elaboração de frases fragmentadas que comprometam a estrutura lógico-gramatical do texto ;  sequência justaposta de ideias sem encaixamentos sintáticos, reproduzindo hábitos da oralidade ;  elaboração de frase com apenas oração subordinada, sem oração principal ;  emprego equivocado do conector ( preposição, conjunção, pronome relativo, alguns advérbios e locuções adverbiais ) comprometendo a expressão da relação lógica entre duas ideias, com prejuízo da clareza do texto ;  emprego do pronome relativo sem a preposição, quando obrigatória ;  repetição ou substituição inadequada de palavras sem o emprego dos recursos oferecidos pela língua ( pronome, advérbio, artigo, sinônimo ) ;  emprego inadequado dos pronomes relativos “ cujo(a) ” e “ onde ” ;  utilização inadequada dos sinais de pontuação, comprometendo a clareza textual ;  inteligibilidade relacionada ao atendimento das exigências de estruturação textual. ( b ) respeito às convenções ortográficas da norma-padrão da Língua Portuguesa – essa competência envolve o domínio das regras de acentuação gráfica e da grafia padrão das palavras ( com ausência de abreviaturas próprias da linguagem da internet ), de acordo com as convenções estabelecidas pela legislação em vigor e consubstanciadas no Vocabulário Ortográfico da Língua Portuguesa, editado pela Academia Brasileira de Letras. ( Vale notar que, nessa edição de 2017, não foi aceita a legislação anterior, no caso das regras relativas ao uso do hífen e da acentuação gráfica ).
Espera -se que o participante :  grafe corretamente as palavras ;  respeite as regras de acentuação gráfica ;  empregue maiúsculas em início de frase, em nomes próprios de pessoas, lugares ou instituições ;  grafe as siglas com letras maiúsculas ou apenas com a primeira letra maiúscula quando formarem uma palavra ;  evite abreviações como p/, vc, tb, pq, tá, né, usadas muitas vezes em escrita informal e na internet ;  obedeça às regras de separação de sílabas no final da linha. ( c ) domínio dos diferentes aspectos morfossintáticos próprios da modalidade escrita formal da norma-padrão da Língua Portuguesa – essa competência envolve : a concordância nominal, a concordância verbal, a regência nominal, a regência verbal, a flexão nominal, a flexão verbal, a correlação entre os tempos verbais, a colocação pronominal e a utilização de sinais de pontuação que contribuam para a organização lógica da frase e do texto.
Espera -se que o participante :  flexione o verbo para estabelecer concordância de número com o sujeito da frase ;  flexione o artigo, o adjetivo e o pronome para concordar em número e em gênero com o substantivo a que se referem ;  observe a regência nominal e a verbal, utilizando a preposição adequada depois de um substantivo, um verbo ou um adjetivo ;  empregue adequadamente o acento grave indicador da crase entre uma preposição e um artigo ( a+a ) ;  obedeça às regras de colocação pronominal ( próclise e ênclise ), distintas dos hábitos da oralidade ou da escrita informal ( exigência de próclise com termo atrator, não exigência de que o pronome oblíquo se ligue ao verbo auxiliar por meio de hífen ) ;  flexione adequadamente verbos, substantivos, adjetivos e pronomes no que diz respeito à expressão das categorias gramaticais ;  flexione os verbos para expressar a correlação de modo e tempo nas estruturas subordinadas.
Com base nesses critérios, foram considerados como desvios de caráter morfossintático, e não como desvios ortográficos, as alterações que envolvem mudança de classe gramatical ou de forma flexional do verbo :  eliminação da marca de infinitivo ( -r- ) e substituição por acento agudo ou ausência total de marca do infinitivo ;  confusão entre “ ão ” e “ am ” nas formas verbais ;  confusão entre “ há ” e “ a ” ;  uso de hífen para separar pronome átono – tanto uso indevido quanto omissão ( exemplo : “ esperasse ”, em lugar de “ espera -se ” ; “ falar -mos ”, no lugar de “ falarmos ” ;  “ esta ” ( no lugar de “ está ” ) ; “ mais ” ( no lugar de “ mas ” ) ; “ e ” ( no lugar de “ é ” ).  verbos “ ter ” e “ vir ” que, na terceira pessoa do plural, não apresentarem o acento circunflexo, serão considerados como desvio de concordância, em aspectos morfossintáticos. ( d ) seleção vocabular adequada à modalidade escrita formal da Língua Portuguesa, exigida pela situação comunicativa – essa competência envolve a precisão na seleção/utilização do vocabulário relacionado à temática solicitada pela questão ; a ausência de marcas da oralidade, como termos de sentido muito genérico ( “ coisa ”, “ negócio ”, “ você ” ) e termos de registros mais informais ( como gírias, jargões, frases feitas, ditados populares, termos regionais ).
Assim, espera -se que o participante respeite a adequação vocabular não usando gírias ou expressões coloquiais, evite repetição desnecessária de palavras e utilize um vocabulário mais formal, como solicitado por um texto dissertativo. \textbf{Observações} : - A inteligibilidade relacionada ao atendimento das exigências de estruturação textual foi avaliada na segunda competência ( relativa aos aspectos textuais ). - Os problemas de coerência textual provocados por uso indevido do vocabulário foram avaliados na terceira competência ( relativa aos aspectos morfossintáticos e vocabulares ). - Não foram considerados como desvios morfossintáticos os problemas de caligrafia ( -a/ -o, -s, -r ). - Não foram considerados os textos grafados integralmente em caixa alta, já que não seria possível distinguir alguma marcação especial para as letras em início de frase. - Cada desvio foi considerado como uma ocorrência, mesmo que dois desvios fossem relativos ao mesmo aspecto linguístico.
A grade de avaliação do desempenho linguístico considerou, portanto, três grandes grupos de competências, segundo os aspectos explicitados anteriormente : a ) domínio das convenções ortográficas : grafia de vogais e consoantes, uso de maiúsculas e minúsculas, emprego do hífen e acentuação gráfica ; b ) domínio dos procedimentos de estruturação textual do ponto de vista microestrutural : organização interna dos períodos, emprego de conectores para a articulação lógica entre os períodos e entre os parágrafos, emprego de marcas de referenciação lexical e pronominal ; utilização dos sinais de pontuação que contribuem para a organização lógica da frase ; c ) domínio das regras de caráter morfossintático estabelecidas como modelares do ponto de vista da modalidade escrita formal da norma-padrão da Língua Portuguesa : concordância nominal e verbal, regência nominal e verbal, colocação pronominal, flexão nominal e verbal, correlação entre tempos e modos verbais, ausência de marcas de oralidade.
A seleção vocabular adequada à modalidade escrita formal da Língua Portuguesa foi incorporada a essa última competência, tendo em vista a intersecção entre as duas do ponto de vista das exigências do registro formal da modalidade escrita da norma-padrão.
Apreciação geral do desempenho dos estudantes : A correção revelou desempenhos distintos dos participantes, marcados pela falta de repertório cultural da maioria dos participantes : - a questão 1 propiciou a oportunidade de maior desenvolvimento do tema solicitado, gerando consequências na estruturação textual, já que os textos foram mais longos, mais elaborados e fluentes.
Entretanto, em virtude da diversidade de abordagens dos textos motivadores ( a falta de penicilina, a recusa no uso da camisinha pelos homens, a maior frequência das mulheres na realização de exames ginecológicos, a transmissão das DST aos bebês durante a gestação ), observaram -se respostas que se configuraram como verdadeiras paráfrases, com comprometimento do caráter \textbf{autoral} desejável em uma questão discursiva de Formação Geral.
Alguns participantes selecionaram um trecho de cada texto motivador para compor suas respostas, outros copiaram trechos inteiros. - a questão 2 teve um comportamento distinto, devido ao caráter polêmico que o caracteriza.
Assim, os textos são mais curtos, com vocabulário repetitivo e preso aos textos motivadores.
Além disso, observou -se grande ocorrência de protesto em função do tema, ora por motivos religiosos ora por preconceito de gênero.
Quanto aos aspectos linguísticos analisados durante esta avaliação, os resultados observados estão descritos nos parágrafos seguintes.
Aspectos ortográficos : O desempenho dos participantes revelou uma diferença muito grande nos dois aspectos analisados nesta competência : baixo índice de desvios da grafia padrão e grande índice de desvios de acentuação.
Em vários casos, ocorre ausência completa de acentuação gráfica.
Os resultados revelam que a tendência dominante entre os universitários brasileiros é a eliminação da acentuação gráfica, talvez motivada pelos hábitos relacionados às redes sociais e pela ausência de esclarecimento dos meios de comunicação, das autoridades e das escolas sobre as decisões do Acordo Ortográfico da Língua Portuguesa de 1990.
Os casos mais sistemáticos de eliminação do acento indicador da sílaba tônica são :  palavras proparoxítonas ( “ sifilis ”, “ proximos ”, “ politicas ”, “ publicas ”, “ transgeneros ” ) ;  palavras paroxítonas terminadas em ditongo crescente ( “ necessario ”, “ noticiarios ”, “ individuo ”, “ dependencia ”, “ varios ”, “ propria ”, “ transmissiveis ”, “ ocorrencia ” ) ;  palavras paroxítonas com hiato ( “ saude ” )  palavras oxítonas ( “ ninguem ”, “ esta ”, “ ate ”, “ tambem ”, “ prevencao ” ).
Por outro lado, destaca -se o uso indevido do acento gráfico em palavras como “ genêro ”, “ melâncolia ”, “ prevênção ”, “ intervênção ”, por exemplo.
Há, também, uma tendência a acentuar paroxítonos como se fossem oxítonos : “ transmissivéis ”.
Quanto ao domínio das convenções relativas à grafia das palavras, observam -se desvios recorrentes, como a hipercorreção pela escolha de “ e ” no lugar de “ i ”, por influência de hábitos da oralidade ( “ descriminação ” no lugar de “ discriminação ”, “ entervenção ” no lugar de “ intervenção ” ) ; desvios de grafia relacionados à variação diastrática podem ser observados em “ estrupo ”, “ subjulgadas ”, “ precoseito ”, “ soubre ”, “ vecendo ”, “ indesencia ”, “ apolojia ”, “ fulga ”, “ dereitos ”, “ sifelis ”.
Observam -se, também, casos de inadequação no uso da maiúscula para destacar determinadas palavras-chave do texto, como “ Transgêneros ”, “ Brasileiros ”, “ Homens ”.
Destaque -se, também, o grande número de participantes que grafam os textos inteiramente em caixa alta.
Vale observar, também, que, ao contrário do que se esperava, não apareceram abreviaturas próprias do “ internetês ”, relacionadas ao uso de redes sociais e e-mails.
Aspectos textuais : Esta competência é a que se revela como a mais problemática entre os participantes, porque são muitos os problemas observados, desvios acumulados durante toda a formação do estudante e que não se resolvem com um estudo autodidata, como acontece com regras ortográficas ou morfossintáticas.
São eles : sequência justaposta de ideias sem encaixamentos sintáticos ; redução drástica de estruturas subordinadas, ao lado do aumento na frequência de estruturas coordenadas e absolutas ; redução no uso de conectores para expressar relações lógicas essenciais à construção do texto, substituídas pela exigência de inferência por parte do interlocutor para suprir a sua ausência ; emprego equivocado de operadores que não estabelecem relações lógicas coerentes entre ideias do texto ; emprego inadequado do pronome relativo ( com omissão da preposição ou a utilização de pronome inadequado, como “ onde ” ) ; repetição exaustiva de termos sem a utilização de procedimentos mais sofisticados de substituição ( hiperonímias, hiponímias, nominalizações, expressões metafóricas ) ; frases fragmentadas que comprometem a estrutura lógico-gramatical ; frases formadas apenas por oração subordinada, sem oração principal.
Um importante aspecto a destacar é o baixo desempenho de uma parte dos participantes em relação à estrutura formal do texto produzido, o que é preocupante ao se levar em conta que são graduandos em fase final de formação.
São frequentes os casos de desvios de estruturação frasal, com uso inadequado ou ausência de conectivos entre parágrafos e entre frases.
Em uma parte dos textos, falta textualidade e domínio do registro padrão da língua.
Na verdade, observam -se relações linguísticas quase agramaticais, como as estabelecidas pela sequência de gerúndios sem o apoio de um ponto de partida para a organização das informações gramaticais e semânticas.
Observou -se que uma parte dos participantes não distribuiu as ideias em parágrafos, talvez devido ao pequeno número de linhas disponibilizadas para a resposta da questão ou, quem sabe, pela suposição de que não seria necessária essa divisão por não se tratar de um texto no modelo de uma redação dissertativo-argumentativa, como solicitado nos vestibulares.
Em função do encaminhamento dos enunciados das duas questões, que solicitaram uma análise do problema e encaminhamentos de políticas públicas, houve uma grande tendência, também, de construção de dois parágrafos desconexos, sem utilização de elementos coesivos adequados à progressão textual.
Quanto à utilização dos mecanismos de referenciação, deve -se destacar a ocorrência de repetições de palavras ou expressões sem a utilização de termos sinônimos ou pronomes, como seria adequado.
Quanto à utilização dos sinais de pontuação, observou -se uma grande precariedade nos textos analisados.
É muito frequente a ocorrência de parágrafos sem marca interna de pontuação para separar os períodos.
Vale observar que não foi penalizada a ausência de vírgula para destacar locuções ou adjuntos adverbiais de pequena extensão deslocados de posição na frase, por ser um uso opcional.
São os seguintes os tipos de problemas encontrados : a ) vírgula : utilização de vírgula para separar o sujeito e o predicado ; ocorrência de apenas uma das vírgulas para separar uma palavra, uma expressão ou uma oração encaixada ; uso de vírgula no lugar do ponto para separar ideias que constituem períodos distintos ; ausência de vírgula para separar elementos de uma enumeração ; ausência de vírgula para separar oração adjetiva explicativa ou utilização inadequada para separar oração adjetiva restritiva ; b ) ponto e vírgula : utilização do ponto e vírgula no lugar de vírgula ; c ) ponto final : ausência de ponto final para separar períodos.
Aspectos morfossintáticos e vocabulares : Em relação à regência, o desvio mais frequente é a falta do sinal indicativo da crase – isso revela que o usuário não tem consciência de que, sob a forma do termo “ a ”, existe a presença de uma contração entre a preposição “ a ” ( exigida pela regência do termo anterior ) e o artigo definido “ a ”.
Um desvio de regência significativo, nos últimos anos, é a utilização inadequada de uma preposição ou sua ausência após o verbo ou o nome ( substantivo ou adjetivo ).
Outro problema relacionado à regência verbal e à nominal, encontrado frequentemente nas questões, foi a ausência de preposição antes de pronome relativo, processo generalizado na modalidade oral da língua, em situações de registro informal.
Apesar da possibilidade de que essa alteração de regência se generalize no padrão escrito da Língua Portuguesa, como já está ocorrendo até em textos jornalísticos, o não emprego da preposição foi considerado inadequado neste processo de avaliação.
A concordância verbal e a concordância nominal apresentam alguns desvios muito frequentes.
Quanto à concordância de número, observou -se ausência de marca ( com sujeito anteposto ou posposto ) ou uso indevido ( uso inadequado da marca de plural comandado pelo núcleo plural da locução adjetiva, apesar de o substantivo que funciona como núcleo do sintagma nominal estar no singular ).
Uma ocorrência que se destacou foi a ausência de acento circunflexo na forma plural do presente do indicativo do verbo “ ter ”, que foi considerada como um desvio na concordância verbal e não na acentuação gráfica.
Quanto à concordância de gênero, vários casos foram observados, normalmente no âmbito de sintagmas nominais longos, em que o adjetivo está afastado do substantivo.
Deve -se destacar o aparecimento da marca de plural em verbos ou adjetivos comandados por adjunto adnominal plural, apesar de serem relacionados a núcleos substantivos no singular, evidenciando um processo de hipercorreção ( exemplo : “ O nome social dos trangêneros podem ser estabelecidos por uma legislação específica ” ).
Quanto à questão da colocação pronominal, foram poucos os desvios observados.
Concluiu -se que, no registro escrito formal, a maioria dos participantes já incorporou regras como a não introdução da frase por um pronome oblíquo e a próclise na presença de um termo atrator.
Não se adotou, entretanto, o padrão excessivamente formal descrito pelas gramáticas normativas em relação à posição do pronome oblíquo em locuções verbais, já que esse uso está muito distante da prática cotidiana, até em textos mais formais.
Quanto aos aspectos vocabulares, alguns tipos de inadequação foram observados : expressões da oralidade ; seleção vocabular incompatível com o contexto, gerando falta de inteligibilidade ; falta de domínio de vocabulário mais abstrato e de maior complexidade, essencial ao desenvolvimento do texto de base dissertativa.
O principal aspecto observado foi a excessiva repetição de certas palavras, revelando limitação de repertório vocabular.
O termo “ pessoa ”, por exemplo, chega a ser repetido até 7 ou 8 vezes em um mesmo texto.
Várias marcas de oralidade foram identificadas, embora não com alta frequência : uso do pronome relativo “ onde ” como relativo universal, falta de artigo definido antes de substantivo, reduções como “ tá ”, “ pra ”, “ pro ”, “ prum ”, expressões informais, eliminação de preposições.
Em função do tema solicitado na questão 1, a sigla DST foi grafada de diferentes maneiras : DST´s, DSTs ou DSTS.
Além disso, houve flutuação de gênero gramatical nos sintagmas que envolviam a sigla : “ os diferentes DSTs ” ou “ as diferentes DSTs ”.
Em função do tema solicitado na questão 2, houve flutuação de gênero gramatical nos sintagmas que envolviam a sua palavra-chave : “ os trans ” ou “ as trans ” ; “ os transgêneros ” ou “ as transgêneras ”.
Todas essas formas foram aceitas, por considerarmos que sua utilização é muito recente e os participantes ainda não sistematizaram essas formas em seu uso da língua. 6.3.2 Componente de Conhecimento Específico Na parte da prova relativa às questões discursivas do Componente de Conhecimento Específico ( Tabela 6.17 ), observa -se que a Média foi bem mais baixa do que para as questões discursivas do Componente de Formação Geral.
Enquanto no Componente de Formação Geral a Média para estudantes de Engenharia de Produção de todo o Brasil foi 56,5, na parte de Conhecimento Específico a Média foi 12,9.
A maior Média deste componente foi obtida pelos estudantes da região Sul ( 14,1 ), e a menor, pelos da região Norte ( 11,4 ).
Quanto à variabilidade das notas, o Desvio padrão de todo o Brasil foi 9,3.
O maior Desvio padrão foi encontrado na região Sul ( 10,0 ), e o menor, nas regiões Norte e Centro-Oeste ( 8,2 ).
A maior nota Máxima, 73,3, foi obtida na região Sudeste.
Nas demais regiões a nota Máxima foi : 48,3 na região Norte, 63,3 na região Nordeste, 66,7 na região Sul e 51,7 na região Centro-Oeste.
A nota Mínima ( 0,0 ) foi obtida por pelo menos um aluno em todas as regiões do Brasil.
A Mediana do Brasil como um todo foi 11,7, a mesma das regiões Nordeste, Sudeste e Centro-Oeste.
A maior Mediana foi obtida na região Sul ( 13,3 ) e a menor na região Norte ( 10,0 ).
O Gráfico 6.10 representa a distribuição das notas nas questões discursivas no Componente de Conhecimento Específico.
A moda desta distribuição ocorre no intervalo [ 0 ; 10 ], com quase 45 \% do total de participantes.
Observa -se que nesse intervalo estão computadas as ocorrências de respostas em branco para as três questões discursivas de Conhecimentos Específicos.
A análise de cada uma destas questões será feita a seguir.
O Gráfico 6.10 representa a distribuição das notas nas questões discursivas no Componente de Conhecimento Específico.
A moda desta distribuição ocorre no intervalo [ 0 ; 10 ], com quase 45 \% do total de participantes.
Observa -se que nesse intervalo estão computadas as ocorrências de respostas em branco para as três questões discursivas de Conhecimentos Específicos.
A análise de cada uma destas questões será feita a seguir. 6.3.2.1 Análise da Questão Discursiva 3 do Componente de Conhecimento Específico Na questão 3, cujos resultados aferidos encontram -se descritos na Tabela 6.18, a Média dos estudantes de todo o Brasil foi 8,7.
Dentre as três discursivas de Conhecimentos Específicos, a questão 3 teve desempenho inferior ao da questão 4 e superior ao da questão 5.
A menor Média nessa questão foi obtida pelos alunos da região Norte ( 6,6 ), enquanto a maior Média foi obtida nas regiões Sudeste e Sul ( 8,9 ).
Quanto à variabilidade das notas, o Desvio padrão de todo o Brasil foi 11,8.
O maior Desvio padrão foi obtido na região Sudeste ( 12,2 ), enquanto o menor foi obtido na região Norte ( 9,4 ).
A nota Máxima, 100,0 pontos, foi alcançada por, pelo menos, um aluno da região Nordeste.
A Mediana do Brasil como um todo foi 5,0, a mesma para todas as regiões.
A nota Mínima ( 0,0 ) também foi a mesma em todas as regiões do Brasil.
O Gráfico 6.11 mostra a distribuição das notas na questão discursiva 3, do Componente de Conhecimento Específico da área de Engenharia de Produção.
Observa -se que em torno de 15 \% dos estudantes deixaram a resposta à questão em branco e cerca de 25 \% receberam nota zero, atingindo quase 40 \% do total de participantes.
Aproximadamente o mesmo percentual de estudantes teve nota no intervalo ( 0 ; 10 ], a moda dessa distribuição. 6.3.2.2 Comentários sobre as respostas à Questão Discursiva 3 A Questão 3 apresentou uma situação típica relacionada à área de Logística e Custos, situação que um Engenheiro de Produção enfrenta na prática de sua profissão, o que reforça a pertinência do tema abordado.
Exigia nível de conhecimento básico desses assuntos, cabendo aos estudantes responder direta e objetivamente ao item ‘ a ’ e modelar o problema de forma a resolver as questões quantitativas dos itens ‘ b ’ e ‘ c ’.
A linguagem do enunciado era clara, e o nível de profundidade exigido pareceu adequado aos formandos.
Todas as informações necessárias para a modelagem do problema para responder aos itens ‘ b ’ e ‘ c ’ estavam disponíveis no enunciado.
Entretanto, o espaço destinado ao desenvolvimento da resposta se mostrou insuficiente, já que ela exigia montagem de equações algébricas.
De fato, o padrão de resposta evidencia a dificuldade em responder à questão no espaço que o estudante tinha disponível.
Pode -se considerar que a questão apresentou nível de dificuldade médio, em comparação com as demais questões de formação específica.
Nessa questão, foi exigido do respondente a interpretação das informações disponibilizadas pelo enunciado e o desenvolvimento de modelos que apoiassem a solução do problema apresentado.
No geral, houve muitas respostas incompletas ou mesmo em branco que apontam para diferentes possibilidades : não entendimento do enunciado, dificuldade em modelar o problema, ou desconhecimento do assunto em pauta.
Em relação às interpretações dadas pelos respondentes, em comparação ao esperado pelo padrão adotado, pode -se tecer as seguintes considerações : - No item ‘ a ’, as respostas seguiram, de maneira geral, o padrão de respostas.
Foi observada, em muitas delas, pouca precisão do que realmente o estudante tencionava destacar – por exemplo, foram citadas vantagens como “ bom para o meio ambiente ”, “ ecologicamente correto ”, “ bom por ser reciclado ”, etc.
Houve também muitos casos de equívoco - muitos afirmaram que os menores custos ( sem especificar quais ) envolvendo a produção de alumínio reciclado seriam uma vantagem desta modalidade de produção.
Também foi frequente a interpretação equivocada dos que trataram a questão como se tivessem sido solicitadas as vantagens do uso de matéria prima primária ( bauxita ), ao invés do uso de matéria prima secundária ( alumínio reciclado ), o que de fato era solicitado na questão.
Percebeu -se, também, a tentativa de responder ao item ‘ a ’ com informações do enunciado da questão.
Por exemplo, muitas respostas apresentaram a existência de dois modais como vantagem ou a existência de custo de transbordo como desvantagem. - No item ‘ b ’, parte considerável dos respondentes conseguiu chegar ao cálculo da produção de alumínio a partir da bauxita - essa parte da resposta exigia uma interpretação mais simples e uma equação elementar.
A partir das respostas dadas, não foi possível diferenciar quem percebeu que os custos de armazenagem e transporte da bauxita eram zero de quem acertou esta parte da questão, entendendo que o custo total era somente o custo de aquisição da matéria prima.
Uma pequena parte dos respondentes acertou o cálculo da aquisição e transporte envolvidos na produção do alumínio reciclado e praticamente nenhum respondente acertou o cálculo da armazenagem - o cálculo deste custo exigia a interpretação do comportamento médio do estoque, considerando lote único de produção.
Também parece não ter ficado claro para os respondentes que o custo de transporte a ser calculado seria o rodoviário, que representava restrição econômica significativa, em comparação aos custos com transporte ferroviário.
A comparação com o transporte ferroviário era objeto do item ‘ c ’.
Por fim, houve muitas respostas que consideraram apenas o custo total unitário ( por tonelada, acrescentado as frações relativas a transporte, armazenagem, etc. ). - No item ‘ c ’, os poucos que desenvolveram a questão conseguiram acertar o cálculo do custo de transporte, mas não chegaram ao cálculo do número de pedidos para se determinar o custo de transbordo no terminal e de armazenagem – para o custo último, cabe a mesma \textbf{observação} feita em relação à interpretação do comportamento do estoque.
Outros respondentes não consideraram todos os itens de custo, tornando o cálculo incompleto.
Finalmente, muitos adotaram a estratégia de calcular o custo unitário, desconsiderando por completo a influência da escala de produção no valor dos custos.
Percebeu -se um grande descolamento entre as respostas dadas pelos respondentes e o que se esperava como resposta padrão nos três itens da questão.
Em síntese, quanto ao item ‘ a ’, ficou evidente que muitas respostas foram dadas sem a devida argumentação, sem justificar, de fato, o que seria considerado vantagem e desvantagem no caso da utilização da matéria-prima secundária.
Nos itens ‘ b ’ e ‘ c ’, que exigiam uso dos dados e equacionamento matemático, o desempenho dos respondentes foi ainda pior.
Notou -se que os respondentes não dominavam conceitos básicos de custos e de logística.
Muitos apresentaram dificuldades de organizar e desenvolver uma linha de raciocínio.
Além do já comentado problema de interpretação do comportamento do estoque, percebeu -se, em muitos casos, ausência de habilidade dos respondentes em organizar dados e informações para estruturar uma solução, mesmo quando se tratava de itens mais simples, como o cálculo do custo da produção de bauxita.
As respostas com notas mais baixas representam os casos em que o estudante deixou a questão em branco, ou respondeu apenas ao item ‘ a ’ e/ou a um dos itens restantes de forma incompleta.
Corresponderam a aproximadamente 90 \% das respostas.
As respostas medianas representam, em geral, os casos em que o estudante respondeu de maneira satisfatória ao item ‘ a ’ e a pelo menos um dos itens restantes ( ‘ b ’ e ‘ c ’ ), apresentando os cálculos desenvolvidos, mesmo que de forma incompleta.
Corresponderam a pouco mais de 5 \% das respostas.
As respostas boas representam os casos em que o estudante respondeu ao item ‘ a ’ de forma satisfatória e desenvolveu os demais itens ( ‘ b ’ e ‘ c ’ ), em parte ou em sua totalidade, de forma correta.
Observou -se grande dificuldade em utilizar os conhecimentos necessários para a solução de problemas da área.
A questão em si não apresentou, de fato, nível elevado de complexidade que justificasse o baixíssimo desempenho dos participantes do exame.
É evidente, por exemplo, a dificuldade enfrentada por eles nas questões referentes à sustentabilidade ambiental, tendo em vista as imprecisões das respostas ao item ‘ a ’.
Em relação aos itens ‘ b ’ e ‘ c ’, mesmo considerando a necessidade de conhecimento de logística e custos, a questão poderia ser, mesmo que parcialmente respondida por pessoas que não fossem da área, bastando, para isso, mostrar habilidades em modelar e resolver problemas, além de conhecimentos básicos em operações algébricas. 6.3.2.3 Análise da Questão Discursiva 4 do Componente de Conhecimento Específico A Tabela 6.19 contém as informações relativas à questão 4 do conjunto de questões do Componente de Conhecimento Específico.
O desempenho dos estudantes nessa questão o melhor dentre as questões discursivas de Conhecimentos Específicos.
A Média geral do Brasil foi 25,2, sendo a menor Média registrada na região Norte ( 23,9 ), e a maior, na região Sul ( 26,5 ).
A nota Máxima ( 100,0 ) foi atingida por, pelo menos, um concluinte das regiões Sudeste e Sul.
A nota Máxima foi 90,0 nas regiões Norte e Nordeste e foi 80,0 na região Centro-Oeste.
A Mediana em todo o Brasil foi 30,0, mesmo valor sendo obtido nas cinco Grandes Regiões.
A nota Mínima de todas as regiões foi zero.
O Gráfico 6.12 representa a distribuição de notas da questão discursiva 4, do Componente de Conhecimento Específico.
Essa distribuição tem moda no intervalo ( 20 ; 30 ], com frequência em torno de 57 \%.
Nessa questão, o percentual de participantes que não responderam ou que receberam nota zero para a resposta dada foi inferior ao da questão 3, ficando em torno de 20 \%. 6.3.2.4 Comentários sobre as respostas à Questão Discursiva 4 O enunciado da Questão 4 apresentou objetivamente o tema abordado - a ergonomia no trabalho.
A linguagem estava clara e precisa, descrevendo o conceito e apresentando uma ilustração com uma situação típica relacionada à ergonomia, que um Engenheiro de Produção enfrenta na prática de sua profissão, o que evidencia a pertinência da questão.
Exigia nível de conhecimento básico de ergonomia, embora o item ‘ b ’ tenha exigido dos respondentes o domínio de alguma abordagem metodológica sobre Análise Ergonômica do Trabalho.
O item ‘ c ’, contudo, referiu -se a um campo de conhecimento correlato à Ergonomia : o Processo de Desenvolvimento do Produto ( PDP ).
Sob este aspecto, houve certa confusão por parte dos respondentes para interpretar a questão.
O texto do enunciado da questão era conceitual, mas a ilustração apresentava uma situação problema específica, o que tornou as perguntas mais independentes do enunciado, em comparação com as demais questões discursivas da prova.
Apesar de as perguntas exigirem respostas objetivas e diretas, o nível de dificuldade, em comparação com as demais questões de formação específica, pode ser considerado de médio para alto.
As muitas \textbf{repostas} incompletas ou mesmo em branco apontam para algumas possibilidades, desde a falta de tempo ao pouco conhecimento sobre o tema.
Em relação às interpretações dadas pelos respondentes, em comparação ao esperado pelo padrão adotado, podem -se tecer as seguintes considerações : - No item ‘ a ’, boa parte dos respondentes fizeram confusão em relação às"queixas físicas"e, diferente do esperado, apontaram problemas relacionados à interação homem-máquina-ambiente.
Por exemplo, citaram sala escura, cadeira baixa, tela do computador afastada ou desalinhada à linha de visada em 90º, dentre outros, e não as consequências desses problemas para o corpo humano ( dores nos membros superiores, dores de cabeça, etc. ), que era o que a pergunta solicitava.
Entretanto, no geral, boa parte dos respondentes pontuou no item. - No item ‘ b ’, quase a totalidade dos respondentes pareceu desconhecer algum método de Avaliação Ergonômica do Trabalho ( AET ).
Boa parte deles utilizou como resposta as medidas cabíveis para melhorar o bem-estar do trabalhador, isto é, somente ações pontuais e não ações como parte de um método estruturado.
Foram observados, assim, exemplos como"ajustar altura da cadeira","apoio para os pés","levantar -se de tempos em tempos", etc., que não representavam a sequência lógica baseada em algum método AET, conforme solicitado na questão. - No item ‘ c ’, nitidamente, os respondentes não conseguiram identificar a relação de PDP com AET.
Muitas respostas citaram ferramentas que não se relacionam com PDP ( a maioria repetiu AET, e outros citaram NIOSH, MOORE, GARG – ferramentas específicas da ergonomia, e não de PDP ) e, ainda, outros citaram ferramentas gerenciais mais genéricas ( como PDCA, SWOT, 5W2H, etc. ).
Em suma, muitos respondentes não perceberam que as ferramentas solicitadas se referiam ao PDP, e não mais à ergonomia.
Pode -se perceber que houve relativo descolamento entre as respostas dadas pelos respondentes e o que se esperava como padrão nos três itens da questão.
Há de se destacar que os itens ‘ b ’ e ‘ c ’ foram os responsáveis pela reduzida média obtida nesta questão : se, por um lado, uma grande parcela dos estudantes respondeu corretamente ( ainda que de forma parcial ) ao item ‘ a ’, a situação se inverteu nos itens ‘ b ’ e ‘ c ’.
Pouquíssimos respondentes responderam adequadamente a esses itens, mesmo de forma parcialmente adequada ( de fato, muitos deixaram em branco um e/ou outro item ).
As respostas que receberam as notas mais baixas – pouco mais do que 80 \% – representam as respostas em branco e aquelas com resposta parcial para o item ‘ a ’ e/ou um dos itens restantes de forma incompleta.
Respostas com notas maiores do que 30,0 representam, em geral, os casos em que o estudante respondeu de maneira satisfatória ao item ‘ a ’ e ao menos parcialmente um ou mais itens restantes ( ‘ b ’ e ‘ c ’ ).
As melhores respostas foram muito poucas e representam os casos em que o estudante respondeu ao item ‘ a ’ de forma satisfatória, mesmo que tenha respondido parcialmente aos dois demais itens ( ‘ b ’ e ‘ c ’ ).
Embora a média desta questão tenha sido ligeiramente superior à da Questão 3, chega -se à conclusão de que os estudantes também apresentaram grande dificuldade em demonstrar os conhecimentos necessários para a solução de problemas da área.
A questão em si não apresentou, de fato, nível elevado de complexidade que justificasse o baixíssimo desempenho desses respondentes.
Exceção pode ser observada no item ‘ b ’ que exigia pleno conhecimento de algum método AET, muito embora tenha sido considerado como resposta correta as que não necessariamente estavam vinculadas a um método específico.
O item ‘ c ’ exigia também conhecimento específico de PDP, mas exigia uma resposta objetiva e direta, bastando ao respondente conhecer ferramentas relacionadas ao Projeto do Produto, um conhecimento fundamental para os engenheiros de produção.
Além disso, no item ‘ c ’, evidenciou -se a dificuldade de interpretação do comando, dado o expressivo número de estudantes que respondeu de forma equivocada. 6.3.2.5 Análise da Questão Discursiva 5 do Componente de Conhecimento Específico A Tabela 6.20 contém as informações relativas à questão 5 do conjunto do Componente de Conhecimento Específico.
O desempenho dos estudantes de todo o Brasil nesta questão foi o mais baixo dentre as três questões discursivas desse componente.
A nota Média dos estudantes de todo o Brasil foi 4,8.
A maior Média foi registrada na região Sul ( 6,8 ), enquanto a menor Média foi registrada na região Norte ( 3,7 ).
Quanto à variabilidade das notas, o Desvio padrão dos alunos do Brasil, como um todo, foi 15,2.
Enquanto o maior desvio foi encontrado na região Sul ( 18,4 ), o menor foi encontrado na região Centro-Oeste ( 12,9 ).
A Mediana para o Brasil e em todas as regiões foi zero, associada ao excessivo número de notas zero, pouco mais de 55 \% das respostas, além dos zeros das respostas em branco, cerca de 32 \%, como mostra o Gráfico 6.13.
Para o conjunto de alunos de Engenharia de Produção do Brasil, a nota Máxima foi 100,0, e a nota Mínima ( 0,0 ).
Os mesmos limites foram encontrados em todas as regiões.
O Gráfico 6.13 representa a distribuição das notas da questão discursiva 5 do Componente de Conhecimento Específico.
Essa distribuição tem moda nas notas zero.
Destaca -se que a frequência de notas zero somada a de estudantes que não responderam à questão chega a quase 90 \% do total de participantes. 6.3.2.6 Comentários sobre as respostas à Questão Discursiva 5 A questão teve enunciado claro, com nível de profundidade adequado ao formando da Área.
Descrevia uma situação típica que o engenheiro de produção enfrenta na prática de sua profissão, especificamente, envolvendo o projeto de um Sistema Produtivo, o que acentua sua pertinência.
O problema fornecia as informações necessárias para que os dois itens propostos pudessem ser respondidos.
Em relação ao item ‘ a ’, o comando era simples e direto, bastando o respondente conhecer o conceito de Tempo de Ciclo para calcular o tempo de ciclo no referido mês de maneira satisfatória.
Em relação ao item ‘ b ’, a questão deixava em aberto os meios pelos quais se desejaria minimizar o desbalanceamento da linha de montagem, havendo, portanto, diversas configurações possíveis e, consequentemente, diversas respostas pertinentes.
Em comparação com as demais questões discursivas de conhecimentos específicos, essa questão foi considerada fácil - não exigia esforço de modelagem e desenvolvimento.
Não obstante, essa relativa facilidade não correspondeu às notas obtidas pelos respondentes.
Essa foi a questão cujas respostas mais diferiram do padrão esperado.
No item ‘ a ’, percebeu -se que grande parte dos respondentes não conhecia o conceito de tempo de ciclo.
Alguns apenas desenharam o fluxo com os tempos envolvidos, sem responder ao que foi solicitado.
Outros calcularam o tempo de ciclo a partir dos tempos fornecidos pelo problema, sem considerar a demanda de 1.200 peças a serem produzidas em 25 dias.
Boa parte somou os tempos de processamento, calculando o tempo em processo.
Houve também casos, em número considerável, de erros na unidade de medida.
Poucos respondentes acertaram, de maneira satisfatória, ao item ‘ a ’.
No item ‘ b ’, observou -se, além de um considerável número de respostas em branco, que os respondentes fizeram confusão sobre o conceito de balanceamento de uma linha de montagem.
Muitos apenas desenharam o fluxo das tarefas e suas relações de precedência com os tempos envolvidos, sem definir as configurações das estações de trabalho agrupadas.
Outros apresentaram opções que não foram consideradas pelo padrão de respostas.
Muitos estudantes, influenciados pelo erro conceitual do item ‘ a ’, os que calcularam o tempo em processo e não o tempo de ciclo, responderam com agrupamentos que \textbf{ultrapassavam} os tempos máximos cabíveis em cada estação de trabalho.
A correção evidenciou a fragilidade conceitual presente na maior parte das respostas.
Muitos apresentaram dificuldades de organizar e desenvolver uma linha de raciocínio, mesmo considerando que as respostas eram simples e diretas.
Além do já comentado problema de interpretação das informações das tarefas e suas precedências, percebeu -se, em muitos casos, ausência de organização de dados e informações.
O mais grave, contudo, foram falhas de interpretação de texto em português - o item ‘ b ’ solicitava um"projeto de centros de trabalho agrupados", o que deveria ter sido atendido com algum tipo de agrupamento, mesmo pelo estudante que, eventualmente, não dominasse o conceito de balanceamento de centro de trabalho.
Importante ressaltar que foi observado nesta questão o maior percentual de respostas em branco ( acima de 30 \% ), talvez por se tratar da última questão discursiva.
O desempenho nesta questão foi baixíssimo.
As respostas que receberam notas abaixo de 30,0 chegaram a quase 95 \%.
As respostas medianas – em torno de 2 \% – foram os casos em que o estudante respondeu de maneira satisfatória ao item ‘ a ’ e pelo menos parcialmente ao item ‘ b ’, apresentando os cálculos desenvolvidos, mesmo que de forma incompleta.
A questão não apresentava, de fato, nível elevado de complexidade que justificasse o baixíssimo desempenho encontrado.
É evidente, por exemplo, a dificuldade na definição do tempo de ciclo, tendo em vista as imprecisões dadas nas respostas do item ‘ a ’, ou mesmo o grande número de respondentes que não pontuaram.
Em relação ao item ‘ b ’, mesmo considerando a necessidade de conhecimento de projeto de sistemas produtivos, em particular, o de balanceamento de linha de produção, a questão poderia ser, ao menos, parcialmente respondida até por pessoas que não fossem da área, bastando, para isso, mostrar habilidades em se agrupar tarefas e propor soluções, mesmo não considerando as restrições do problema. 6.3.3 Considerações Finais Como considerações finais referentes ao processo de avaliação do desempenho de estudantes de Engenharia de Produção no presente exame, três aspectos devem ser realçados.
O primeiro está relacionado ao propósito dessa avaliação.
Sabe -se que os alunos ainda veem a avaliação do seu curso pelo Enade como um processo burocrático, e grande parte deles ainda não se conscientizou da sua importância para aprimoramento da qualidade dos cursos de ensino superior.
Tal conscientização vem melhorando ao longo dos últimos anos, porém é comum encontrar no meio discente e docente opiniões contrárias ao processo.
Muitas questões em branco e o baixo rendimento nas questões discursivas são prováveis evidências deste obstáculo a ser vencido.
Este aspecto remete ao segundo aspecto relevante : o papel que as diferentes Instituições de Ensino Superior ( IES ) têm desempenhado na formação dos engenheiros de produção de nosso país.
Muitas das lacunas apresentadas pelos estudantes indicam que os cursos de Engenharia de Produção devem se estruturar melhor para oferecer conteúdo aderente às necessidades do mercado, sem descolar, contudo, dos avanços de conhecimento observados em nível mundial e registrados por meio de artigos científicos.
Sobretudo, as IES precisam apoiar o processo de conscientização da importância da avaliação do ensino superior e preparar melhor seus alunos para realização de exames que visam a avaliar suas competências e os conhecimentos adquiridos.
Não se defende, contudo, um"preparatório para o ENADE"puro e simples, já que se espera que a avaliação do Enade seja orgânica ao currículo dos cursos.
Os Coordenadores Acadêmicos das IES precisam se empenhar nos processos de revisão dos currículos, sem deixar de reforçar o conteúdo considerado clássico da Engenharia de Produção.
Ficou evidente nesta avaliação que nossos formandos não dominam problemas básicos, como os relacionados a custos, logística, ergonomia e sistemas de produção.
Além do reforço nas disciplinas consideradas clássicas, conteúdos decorrentes da aplicação de tais conceitos, como meio ambiente e sustentabilidade, empreendedorismo, etc. não podem ser desconsiderados.
Enfim, os Coordenadores Acadêmicos possuem um desafio importante que deve apontar para o aprimoramento dos conteúdos centrais da Engenharia de Produção, ao mesmo tempo em que articular tais conteúdos com assuntos contemporâneos, inerentes à própria demanda que a sociedade nos coloca.
Por fim, o terceiro aspecto a ser ressaltado remete à qualidade da Educação Básica oferecida no país.
É evidente que muitas das dificuldades observadas na resolução das questões são fruto de limitações na interpretação de textos e na capacidade de expressão escrita.
Em muitos casos, respostas satisfatórias poderiam ter sido dadas mesmo por um não especialista no assunto.
Os erros cometidos pelos respondentes foram considerados, na grande maioria, grosseiros e inadmissíveis para um recém-formado em Engenharia.
Uma simples interpretação de texto seria suficiente para responder boa parte da prova discursiva, o que expõe as fragilidades do nosso sistema de ensino trazidas pelo estudante desde a educação básica até o ensino superior.

% matched lemmas: advogado, assento, autoral, cartório, casal, complicador, observação, repor, ultrapassar
\end{textsample}
